\documentclass[11pt,paper=a4,answers]{exam}
\usepackage{graphicx,lastpage}
\usepackage{upgreek}
\usepackage{censor}
\usepackage{gensymb}
\usepackage{amsmath}
\usepackage{multicol}
\usepackage{enumitem}
\usepackage{tasks}
\usepackage{adjustbox}
\usepackage{xcolor}
\usepackage{tfrupee}
\setlength{\voffset}{0.25in}
\usepackage{tabularx}
\newcolumntype{Y}{>{\centering\arraybackslash}X}
%==============================================================
% Customise Non-Essential Packages Below
% =============================================================
\usepackage[most]{tcolorbox}
\usepackage{eso-pic}
\usepackage{tikz}
\AddToShipoutPictureBG{%
\begin{tikzpicture}[remember picture,overlay]
\foreach \x in {-8,-4,0,4,8,12,16,20}
\foreach \y in {-8,-4,0,4,8,12,16,20} {
\node[opacity=0.05,rotate=45,anchor=center] at ([xshift=\x cm,yshift=\y cm]current page.center) {%
\begin{minipage}{2cm}
\centering
\includegraphics[width=2cm]{images/logo_black.png}\\
\small preppeo.com
\end{minipage}%
};
}
\end{tikzpicture}%
}
\printanswers

%==============================================================
% Customise Non-Essential Packages Above
% =============================================================
\definecolor{svgray}{rgb}{0.90, 0.90, 0.90}
\graphicspath{{images/}}

\usepackage[normalem]{ulem}
\renewcommand\ULthickness{1pt}   %%---> For changing thickness of underline
\setlength\ULdepth{3.5ex}%\maxdimen ---> For changing depth of underline
\renewcommand{\baselinestretch}{1}
\chead{
\includegraphics[scale=0.1,height=2cm ]{logo.png}
}
\pagestyle{empty}

\pagestyle{headandfoot}
\headrule
\cfoot{W: https://preppeo.com; M: +91-8130711689}

\begin{document}
%==============================================================
% Customise Question paper details below this
% =============================================================
\begin{center}
    \centering
    \renewcommand{\arraystretch}{1.5}
    \begin{tabularx}{\textwidth}{YYY}
    & \normalsize\bf{{\Large{{Quant}}}} & \\
    By Preppeo & \normalsize\bf{{sat\_lid\_013}} & SAT\\
    \normalsize{{\textbf{{Unit:}} Algebra}} & \normalsize{{\textbf{{Chapter:}} Linear functions: slope, intercepts, graphs}} & \normalsize{{\textbf{{Topic:Linear Growth \& Decay Models}} }} \\
    \end{tabularx}
\end{center}
\par\noindent\rule{\textwidth}{1pt}
% =============================================================
% Customise Question paper details above this
% =============================================================

\begin{questions}
\pointsinrightmargin
\pointsdroppedatright
\marksnotpoints
\pointpoints{mark}{marks}
\pointformat{\boldmath\themarginpoints}
\bracketedpoints

% =============================================================
% Question Begin
% =============================================================


% --- PART 1: Core Growth and Decay Interpretation (1-25) ---

% Question 1
% Category: Practice | Difficulty: Medium | Question Type: Multiple Choice
\question
A scientist is tracking the population of a specific type of bacteria in a controlled environment. The population $P$ can be modeled by the equation $P = 150t + 2,000$, where $t$ is the number of hours since the start of the observation. Which of the following is the best interpretation of the number 150 in this context?
\begin{tasks}(1)
    \task The initial number of bacteria when the observation began.
    \task The number of hours it takes for the population to reach 2,000.
    \task The increase in the number of bacteria each hour.
    \task The total number of bacteria after 150 hours.
\end{tasks}

\begin{solution}
\textbf{Conceptual Explanation:} \\
In a linear model $y = mx + b$, the slope $m$ represents the rate of change—how much the dependent variable ($P$) changes for every one unit of the independent variable ($t$).

\vspace{10pt}
\textbf{Calculation and Logic:} \\
$P = 150t + 2,000$. \\
The coefficient of $t$ is 150. Since $t$ is measured in hours, 150 represents the rate of increase per hour.

\vspace{15pt}
Answer: \colorbox{yellow!100}{C}
\end{solution}

\vspace{20pt}

% Question 2
% Category: Practice | Difficulty: Medium | Question Type: Multiple Choice
\question
A car’s value depreciates linearly over time. The value $V$, in dollars, of a certain car $t$ years after it was purchased is given by $V = 28,500 - 2,100t$. What is the best interpretation of 28,500 in this context?
\begin{tasks}(1)
    \task The value of the car after 1 year.
    \task The amount the car's value decreases each year.
    \task The number of years until the car has no value.
    \task The value of the car at the time of purchase.
\end{tasks}

\begin{solution}
\textbf{Conceptual Explanation:} \\
The $y$-intercept ($b$) represents the "starting" or "initial" value when the time ($t$) is 0.

\vspace{10pt}
\textbf{Calculation and Logic:} \\
At the time of purchase, $t = 0$. \\
$V = 28,500 - 2,100(0) = 28,500$. \\
This is the original price of the car.

\vspace{15pt}
Answer: \colorbox{yellow!100}{D}
\end{solution}

\vspace{20pt}

% Question 3
% Category: Practice | Difficulty: Hard | Question Type: Multiple Choice
\question
A water tank holds 1,200 gallons of water and is being drained at a constant rate. After 5 hours, the tank contains 950 gallons. Which of the following linear functions represents the amount of water $W$, in gallons, remaining in the tank after $t$ hours?
\begin{tasks}(2)
    \task $W = 1,200 - 50t$
    \task $W = 1,200 - 250t$
    \task $W = 950 - 50t$
    \task $W = 1,200 - 5t$
\end{tasks}

\begin{solution}
\textbf{Conceptual Explanation:} \\
First, find the rate of change (slope). The rate is the change in water divided by the change in time. Since it is being drained, the rate will be negative.

\vspace{10pt}
\textbf{Calculation and Logic:} \\
Change in water = $950 - 1,200 = -250$ gallons. \\
Change in time = 5 hours. \\
Rate = $-250 / 5 = -50$ gallons per hour. \\
Initial amount ($b$) = 1,200. \\
Equation: $W = 1,200 - 50t$.

\vspace{15pt}
Answer: \colorbox{yellow!100}{A}
\end{solution}

\vspace{20pt}

% Question 4
% Category: Practice | Difficulty: Easy | Question Type: Multiple Choice
\question
A construction worker’s total daily pay $P$, in dollars, is given by $P = 25h + 40$, where $h$ is the number of hours worked. The constant 40 represents a travel stipend. If the worker’s hourly rate increases by \$5, what will be the new equation for the worker’s daily pay?
\begin{tasks}(2)
    \task $P = 25h + 45$
    \task $P = 30h + 40$
    \task $P = 30h + 45$
    \task $P = 5h + 40$
\end{tasks}

\begin{solution}
\textbf{Conceptual Explanation:} \\
The hourly rate is the slope of the linear function. Increasing the "rate" means adding to the coefficient of $h$.

\vspace{10pt}
\textbf{Calculation and Logic:} \\
Original rate = 25. \\
New rate = $25 + 5 = 30$. \\
The fixed stipend (y-intercept) remains 40. \\
New Equation: $P = 30h + 40$.

\vspace{15pt}
Answer: \colorbox{yellow!100}{B}
\end{solution}

\vspace{20pt}

% Question 5
% Category: Practice | Difficulty: Hard | Question Type: Multiple Choice
\question

The graph above shows the temperature $T$, in degrees Celsius, of a heated liquid as it cools over time $m$, in minutes. Based on the graph, which of the following is true?
\begin{tasks}(1)
    \task The liquid cools at a rate of 20 degrees per minute.
    \task The initial temperature was 40 degrees Celsius.
    \task The liquid cools at a rate of 5 degrees per minute.
    \task It will take 10 minutes for the liquid to reach 0 degrees.
\end{tasks}

\begin{solution}
\textbf{Conceptual Explanation:} \\
Determine the slope from the graph. The slope represents the cooling rate.

\vspace{10pt}
\textbf{Calculation and Logic:} \\
Initial point: $(0, 60)$. After 4 minutes: $(4, 40)$. \\
Change in Temp = $40 - 60 = -20$. \\
Change in time = 4 minutes. \\
Rate = $-20 / 4 = -5$ degrees per minute.

\vspace{15pt}
Answer: \colorbox{yellow!100}{C}
\end{solution}

\vspace{20pt}

% Question 6
% Category: Practice | Difficulty: Medium | Question Type: Multiple Choice
\question
A hiker is climbing a mountain at a constant rate. Their elevation $E$, in feet, $h$ hours after they started is given by $E = 800h + 1,200$. What was the hiker’s elevation when they started?
\begin{tasks}(4)
    \task 0 feet
    \task 800 feet
    \task 1,200 feet
    \task 2,000 feet
\end{tasks}

\begin{solution}
\textbf{Conceptual Explanation:} \\
"When they started" corresponds to $h = 0$.

\vspace{10pt}
\textbf{Calculation and Logic:} \\
$E = 800(0) + 1,200 = 1,200$.

\vspace{15pt}
Answer: \colorbox{yellow!100}{C}
\end{solution}

\vspace{20pt}

% Question 7
% Category: Practice | Difficulty: Hard | Question Type: Multiple Choice
\question
A library has 5,000 books and adds $n$ new books each month. After 12 months, the library has 5,720 books. What is the value of $n$?
\begin{tasks}(4)
    \task 60
    \task 72
    \task 120
    \task 572
\end{tasks}

\begin{solution}
\textbf{Conceptual Explanation:} \\
The total increase over a period of time is the rate multiplied by the time. Set up a linear equation to solve for the rate $n$.

\vspace{10pt}
\textbf{Calculation and Logic:} \\
Total books = Initial + (Rate $\times$ Months) \\
$5,720 = 5,000 + n(12)$ \\
$720 = 12n$ \\
$n = 60$

\vspace{15pt}
Answer: \colorbox{yellow!100}{A}
\end{solution}

\vspace{20pt}

% Question 8
% Category: Practice | Difficulty: Medium | Question Type: Multiple Choice
\question
A smartphone battery charges at a constant rate. The battery percentage $B$ after $m$ minutes of charging is given by $B = 0.8m + 15$. How many minutes will it take for the battery to be fully charged ($100\%$)?
\begin{tasks}(4)
    \task 85 minutes
    \task 106.25 minutes
    \task 125 minutes
    \task 15 minutes
\end{tasks}

\begin{solution}
\textbf{Conceptual Explanation:} \\
Substitute the target value for $B$ and solve for the time $m$.

\vspace{10pt}
\textbf{Calculation and Logic:} \\
$100 = 0.8m + 15$ \\
$85 = 0.8m$ \\
$m = 85 / 0.8 = 106.25$

\vspace{15pt}
Answer: \colorbox{yellow!100}{B}
\end{solution}

\vspace{20pt}

% Question 9
% Category: Practice | Difficulty: Hard | Question Type: Multiple Choice
\question

The graph above shows the height $H$, in feet, of two different trees over $t$ years. Tree A is modeled by $A(t) = 1.5t + 10$ and Tree B is modeled by $B(t) = 2t + 4$. After how many years will both trees be the same height?
\begin{tasks}(4)
    \task 6 years
    \task 10 years
    \task 12 years
    \task 14 years
\end{tasks}

\begin{solution}
\textbf{Conceptual Explanation:} \\
When two linear models are "the same," their equations are equal. Set the two functions equal to each other and solve for $t$.

\vspace{10pt}
\textbf{Calculation and Logic:} \\
$1.5t + 10 = 2t + 4$ \\
$6 = 0.5t$ \\
$t = 12$

\vspace{15pt}
Answer: \colorbox{yellow!100}{C}
\end{solution}

\vspace{20pt}

% Question 10
% Category: Practice | Difficulty: Medium | Question Type: Multiple Choice
\question
A candle is 12 inches tall and burns at a rate of 0.75 inches per hour. Which inequality represents the number of hours $h$ the candle can burn such that its height is at least 3 inches?
\begin{tasks}(2)
    \task $12 - 0.75h \ge 3$
    \task $12 + 0.75h \le 3$
    \task $0.75h \ge 3$
    \task $12 - 0.75h \le 3$
\end{tasks}

\begin{solution}
\textbf{Conceptual Explanation:} \\
The current height is the initial height minus the amount burned. "At least" translates to $\ge$.

\vspace{10pt}
\textbf{Calculation and Logic:} \\
Initial = 12. Amount burned = $0.75 \times h$. \\
Height = $12 - 0.75h$. \\
Condition: $12 - 0.75h \ge 3$.

\vspace{15pt}
Answer: \colorbox{yellow!100}{A}
\end{solution}

\vspace{20pt}

% Question 11
% Category: Practice | Difficulty: Hard | Question Type: Multiple Choice
\question
A salesperson earns a base salary plus a commission for every car sold. The linear function $S(c) = 500c + 25,000$ represents the salesperson’s annual salary $S$, in dollars, for selling $c$ cars. If the salesperson wants to increase their annual salary by \$5,000, how many more cars must they sell?
\begin{tasks}(4)
    \task 10
    \task 50
    \task 5
    \task 60
\end{tasks}

\begin{solution}
\textbf{Conceptual Explanation:} \\
The slope (500) tells us how much the salary increases per car. Divide the target increase by the rate per car.

\vspace{10pt}
\textbf{Calculation and Logic:} \\
Salary increase = 5,000. \\
Rate = 500 per car. \\
Cars needed = $5,000 / 500 = 10$.

\vspace{15pt}
Answer: \colorbox{yellow!100}{A}
\end{solution}

\vspace{20pt}

% Question 12
% Category: Practice | Difficulty: Medium | Question Type: Multiple Choice
\question
A pool is being filled with water. The volume of water $V$, in gallons, after $m$ minutes is $V = 12m + 200$. What does the 200 represent in this context?
\begin{tasks}(1)
    \task The number of gallons added to the pool each minute.
    \task The time it takes for the pool to be full.
    \task The amount of water already in the pool when the filling started.
    \task The maximum capacity of the pool.
\end{tasks}

\begin{solution}
\textbf{Conceptual Explanation:} \\
The constant term in a linear growth model is the initial value at time zero.

\vspace{10pt}
\textbf{Calculation and Logic:} \\
When $m = 0$, $V = 200$. This is the starting volume.

\vspace{15pt}
Answer: \colorbox{yellow!100}{C}
\end{solution}

\vspace{20pt}

% Question 13
% Category: Practice | Difficulty: Hard | Question Type: Multiple Choice
\question
A printer has 400 sheets of paper and uses them at a rate of 8 sheets per minute. After $x$ minutes, the printer has fewer than 100 sheets left. Which inequality represents this situation?
\begin{tasks}(2)
    \task $400 - 8x < 100$
    \task $400 - 8x > 100$
    \task $8x < 300$
    \task $400 + 8x < 100$
\end{tasks}

\begin{solution}
\textbf{Conceptual Explanation:} \\
The remaining amount is initial minus (rate $\times$ time). "Fewer than" means $<$.

\vspace{10pt}
\textbf{Calculation and Logic:} \\
Initial = 400. Usage = $8x$. \\
Remaining = $400 - 8x$. \\
Inequality: $400 - 8x < 100$.

\vspace{15pt}
Answer: \colorbox{yellow!100}{A}
\end{solution}

\vspace{20pt}

% Question 14
% Category: Practice | Difficulty: Medium | Question Type: Multiple Choice
\question

The graph above shows the balance of a gift card $B$ after $n$ coffee purchases. What is the cost of each coffee?
\begin{tasks}(4)
    \task \$100
    \task \$10
    \task \$1
    \task \$5
\end{tasks}

\begin{solution}
\textbf{Conceptual Explanation:} \\
The cost of each coffee is the slope (the rate at which the balance decreases per purchase).

\vspace{10pt}
\textbf{Calculation and Logic:} \\
Total drop = \$100. \\
Total purchases to reach zero = 10. \\
Cost per purchase = $100 / 10 = 10$.

\vspace{15pt}
Answer: \colorbox{yellow!100}{B}
\end{solution}

\vspace{20pt}

% Question 15
% Category: Practice | Difficulty: Hard | Question Type: Multiple Choice
\question
A subscription service costs \$15 per month plus a one-time activation fee of \$30. Which function gives the total cost $C$ for $m$ months of service?
\begin{tasks}(2)
    \task $C(m) = 30m + 15$
    \task $C(m) = 15m + 30$
    \task $C(m) = 45m$
    \task $C(m) = 15(m + 30)$
\end{tasks}

\begin{solution}
\textbf{Conceptual Explanation:} \\
The "per month" cost is the rate (slope), and the "one-time" fee is the initial value (y-intercept).

\vspace{10pt}
\textbf{Calculation and Logic:} \\
Rate = 15. Initial = 30. \\
Equation: $C = 15m + 30$.

\vspace{15pt}
Answer: \colorbox{yellow!100}{B}
\end{solution}

\vspace{20pt}

% Question 16
% Category: Practice | Difficulty: Medium | Question Type: Multiple Choice
\question
The height of a candle $h$, in centimeters, after burning for $t$ minutes is $h = -0.2t + 25$. How many centimeters of the candle burn every 10 minutes?
\begin{tasks}(4)
    \task 0.2 cm
    \task 2 cm
    \task 25 cm
    \task 23 cm
\end{tasks}

\begin{solution}
\textbf{Conceptual Explanation:} \\
The slope represents the change per one minute. To find the change for 10 minutes, multiply the slope by 10.

\vspace{10pt}
\textbf{Calculation and Logic:} \\
Rate = 0.2 cm per minute. \\
For 10 minutes = $0.2 \times 10 = 2$ cm.

\vspace{15pt}
Answer: \colorbox{yellow!100}{B}
\end{solution}

\vspace{20pt}

% Question 17
% Category: Practice | Difficulty: Hard | Question Type: Multiple Choice
\question
A car travels at a constant speed of 65 miles per hour. The distance $D$ from its destination after $t$ hours is $D = 325 - 65t$. What does the 325 represent?
\begin{tasks}(1)
    \task The total time of the trip.
    \task The speed of the car.
    \task The total distance of the trip when $t=0$.
    \task The distance traveled after 1 hour.
\end{tasks}

\begin{solution}
\textbf{Conceptual Explanation:} \\
In a decay or distance-remaining model, the constant term is the initial distance from the goal.

\vspace{10pt}
\textbf{Calculation and Logic:} \\
When $t = 0$, $D = 325$. This is the starting distance.

\vspace{15pt}
Answer: \colorbox{yellow!100}{C}
\end{solution}

\vspace{20pt}

% Question 18
% Category: Practice | Difficulty: Medium | Question Type: Multiple Choice
\question
An athlete's weight $W$, in pounds, after $x$ weeks of a training program is $W = 195 - 1.5x$. How many weeks will it take the athlete to reach 180 pounds?
\begin{tasks}(4)
    \task 10 weeks
    \task 15 weeks
    \task 1.5 weeks
    \task 12 weeks
\end{tasks}

\begin{solution}
\textbf{Conceptual Explanation:} \\
Substitute the target weight for $W$ and solve for $x$.

\vspace{10pt}
\textbf{Calculation and Logic:} \\
$180 = 195 - 1.5x$ \\
$-15 = -1.5x$ \\
$x = 10$

\vspace{15pt}
Answer: \colorbox{yellow!100}{A}
\end{solution}

\vspace{20pt}

% Question 19
% Category: Practice | Difficulty: Hard | Question Type: Multiple Choice
\question
A business's profit $P$ increases by \$2,000 for every additional 100 units sold. If the profit is \$10,000 when 400 units are sold, which function represents profit as a function of units $u$?
\begin{tasks}(2)
    \task $P(u) = 20u + 2,000$
    \task $P(u) = 200u + 10,000$
    \task $P(u) = 20u + 10,000$
    \task $P(u) = 20u + 8,000$
\end{tasks}

\begin{solution}
\textbf{Conceptual Explanation:} \\
First find the rate per unit (slope), then use the point (400, 10,000) to find the y-intercept.

\vspace{10pt}
\textbf{Calculation and Logic:} \\
Rate = $2,000 / 100 = 20$ per unit. \\
$10,000 = 20(400) + b$ \\
$10,000 = 8,000 + b \Rightarrow b = 2,000$. \\
Function: $P(u) = 20u + 2,000$.

\vspace{15pt}
Answer: \colorbox{yellow!100}{A}
\end{solution}

\vspace{20pt}

% Question 20
% Category: Practice | Difficulty: Medium | Question Type: Multiple Choice
\question
A plane is descending at a constant rate. Its altitude $A$, in feet, $t$ minutes after beginning its descent is $A = 30,000 - 1,500t$. How many minutes after beginning its descent will the plane be at an altitude of 15,000 feet?
\begin{tasks}(4)
    \task 10 minutes
    \task 15 minutes
    \task 20 minutes
    \task 5 minutes
\end{tasks}

\begin{solution}
\textbf{Conceptual Explanation:} \\
Set the altitude $A$ to 15,000 and solve for the time $t$.

\vspace{10pt}
\textbf{Calculation and Logic:} \\
$15,000 = 30,000 - 1,500t$ \\
$-15,000 = -1,500t$ \\
$t = 10$

\vspace{15pt}
Answer: \colorbox{yellow!100}{A}
\end{solution}

\vspace{20pt}

% Question 21
% Category: Practice | Difficulty: Hard | Question Type: Multiple Choice
\question
The cost of renting a truck is \$40 per day plus \$0.50 per mile. If a person rents a truck for one day and drives $d$ miles, their total cost is $C$. If the per-mile rate increases by 20\%, what is the new cost equation?
\begin{tasks}(2)
    \task $C = 40 + 0.60d$
    \task $C = 48 + 0.50d$
    \task $C = 40 + 0.70d$
    \task $C = 48 + 0.60d$
\end{tasks}

\begin{solution}
\textbf{Conceptual Explanation:} \\
The per-mile rate is the slope. Calculate the new rate by increasing the original slope by 20\%.

\vspace{10pt}
\textbf{Calculation and Logic:} \\
Original slope = 0.50. \\
Increase = $0.50 \times 0.20 = 0.10$. \\
New slope = $0.50 + 0.10 = 0.60$. \\
Equation: $C = 40 + 0.60d$.

\vspace{15pt}
Answer: \colorbox{yellow!100}{A}
\end{solution}

\vspace{20pt}

% Question 22
% Category: Practice | Difficulty: Medium | Question Type: Multiple Choice
\question
A savings account starts with \$500 and \$25 is deposited every week. Which of the following represents the total amount $S$ in the account after $w$ weeks?
\begin{tasks}(2)
    \task $S = 500w + 25$
    \task $S = 25w + 500$
    \task $S = 525w$
    \task $S = 25(w + 500)$
\end{tasks}

\begin{solution}
\textbf{Conceptual Explanation:} \\
The starting amount is the $y$-intercept, and the weekly deposit is the slope.

\vspace{10pt}
\textbf{Calculation and Logic:} \\
$b = 500, m = 25$. \\
Equation: $S = 25w + 500$.

\vspace{15pt}
Answer: \colorbox{yellow!100}{B}
\end{solution}

\vspace{20pt}

% Question 23
% Category: Practice | Difficulty: Hard | Question Type: Multiple Choice
\question
A pool contains 8,000 gallons of water and is being filled at a rate of 15 gallons per minute. After how many hours will the pool contain 10,700 gallons?
\begin{tasks}(4)
    \task 3 hours
    \task 180 hours
    \task 2.5 hours
    \task 45 hours
\end{tasks}

\begin{solution}
\textbf{Conceptual Explanation:} \\
Find the total time in minutes first, then convert that time into hours.

\vspace{10pt}
\textbf{Calculation and Logic:} \\
Gallons to add = $10,700 - 8,000 = 2,700$. \\
Minutes = $2,700 / 15 = 180$ minutes. \\
Hours = $180 / 60 = 3$ hours.

\vspace{15pt}
Answer: \colorbox{yellow!100}{A}
\end{solution}

\vspace{20pt}

% Question 24
% Category: Practice | Difficulty: Medium | Question Type: Multiple Choice
\question
A person is paying off a \$2,400 loan by making monthly payments of \$150. Which function represents the remaining balance $B$ after $m$ months?
\begin{tasks}(2)
    \task $B(m) = 2,400 + 150m$
    \task $B(m) = 150m - 2,400$
    \task $B(m) = 2,400 - 150m$
    \task $B(m) = 2,250m$
\end{tasks}

\begin{solution}
\textbf{Conceptual Explanation:} \\
The balance starts at the loan amount and decreases by the payment amount each month.

\vspace{10pt}
\textbf{Calculation and Logic:} \\
Initial = 2,400. Decrease = $150 \times m$. \\
Equation: $B = 2,400 - 150m$.

\vspace{15pt}
Answer: \colorbox{yellow!100}{C}
\end{solution}

\vspace{20pt}

% Question 25
% Category: Practice | Difficulty: Hard | Question Type: Multiple Choice
\question
A linear model $y = mx + b$ is used to predict the value of a stock. If the stock's value increases by \$4 every 5 days and was worth \$50 on Day 0, what is the value of $m$?
\begin{tasks}(4)
    \task 4
    \task 5
    \task 0.8
    \task 1.25
\end{tasks}

\begin{solution}
\textbf{Conceptual Explanation:} \\
The slope $m$ is the change in the output ($y$) per one unit of the input ($x$).

\vspace{10pt}
\textbf{Calculation and Logic:} \\
Change in $y = 4$. Change in $x = 5$. \\
$m = 4 / 5 = 0.8$.

\vspace{15pt}
Answer: \colorbox{yellow!100}{C}
\end{solution}


% --- PART 2: Comparative Models and Variable Analysis (26-50) ---

% Question 26
% Category: Practice | Difficulty: Medium | Question Type: Multiple Choice
\question
A delivery service uses the model $C = 2.5d + 7.5$ to calculate the delivery cost $C$, in dollars, for a distance of $d$ miles. Which of the following is a true statement based on the model?
\begin{tasks}(1)
    \task The cost increases by \$7.50 for every mile traveled.
    \task For every 2-mile increase in distance, the cost increases by \$5.00.
    \task The distance must be at least 7.5 miles.
    \task A delivery of 0 miles would cost \$2.50.
\end{tasks}

\begin{solution}
\textbf{Conceptual Explanation:} \\
The slope ($2.5$) represents the cost per mile. To check a 2-mile increase, multiply the slope by 2.

\vspace{10pt}
\textbf{Calculation and Logic:} \\
Slope = 2.50 per mile. \\
2 miles $\times$ 2.50 = 5.00. \\
Therefore, the cost increases by \$5.00 for every 2 miles.

\vspace{15pt}
Answer: \colorbox{yellow!100}{B}
\end{solution}

\vspace{20pt}

% Question 27
% Category: Practice | Difficulty: Hard | Question Type: Multiple Choice
\question
A computer's value $V$ decreases by \$150 each year. After 3 years, the computer is worth \$550. Which equation represents the value of the computer $t$ years after it was purchased?
\begin{tasks}(2)
    \task $V = 150t + 550$
    \task $V = -150t + 1,000$
    \task $V = -150t + 550$
    \task $V = -150t + 400$
\end{tasks}

\begin{solution}
\textbf{Conceptual Explanation:} \\
Use the rate of change and the given point (3, 550) to find the initial value ($y$-intercept).

\vspace{10pt}
\textbf{Calculation and Logic:} \\
$m = -150$. \\
$550 = -150(3) + b$ \\
$550 = -450 + b \Rightarrow b = 1,000$. \\
Equation: $V = -150t + 1,000$.

\vspace{15pt}
Answer: \colorbox{yellow!100}{B}
\end{solution}

\vspace{20pt}

% Question 28
% Category: Practice | Difficulty: Medium | Question Type: Multiple Choice
\question
The remaining length $L$ of a yarn roll, in meters, after knitting $s$ sweaters is $L = 500 - 45s$. If a knitter has 140 meters of yarn left, how many sweaters did they knit?
\begin{tasks}(4)
    \task 6
    \task 8
    \task 10
    \task 12
\end{tasks}

\begin{solution}
\textbf{Conceptual Explanation:} \\
Substitute the remaining length for $L$ and solve for $s$.

\vspace{10pt}
\textbf{Calculation and Logic:} \\
$140 = 500 - 45s$ \\
$-360 = -45s$ \\
$s = 8$.

\vspace{15pt}
Answer: \colorbox{yellow!100}{B}
\end{solution}

\vspace{20pt}

% Question 29
% Category: Practice | Difficulty: Hard | Question Type: Multiple Choice
\question
Two companies offer internet plans. Plan A costs \$50 per month. Plan B costs \$20 per month plus \$0.10 for every gigabyte $g$ of data used. For what amount of data is Plan B more expensive than Plan A?
\begin{tasks}(4)
    \task $g > 300$
    \task $g < 300$
    \task $g > 700$
    \task $g < 700$
\end{tasks}

\begin{solution}
\textbf{Conceptual Explanation:} \\
Set up an inequality where the expression for Plan B is greater than Plan A.

\vspace{10pt}
\textbf{Calculation and Logic:} \\
Plan B: $20 + 0.10g$. Plan A: 50. \\
$20 + 0.10g > 50$ \\
$0.10g > 30$ \\
$g > 300$.

\vspace{15pt}
Answer: \colorbox{yellow!100}{A}
\end{solution}

\vspace{20pt}

% Question 30
% Category: Practice | Difficulty: Medium | Question Type: Multiple Choice
\question
The height of a stack of $n$ identical books is $H = 2.5n + 1$. What does the 1 represent in this context?
\begin{tasks}(1)
    \task The thickness of each book.
    \task The total height of the stack.
    \task The thickness of the bottom cover or a base.
    \task The weight of one book.
\end{tasks}

\begin{solution}
\textbf{Conceptual Explanation:} \\
The $y$-intercept represents the height when the number of books is zero.

\vspace{10pt}
\textbf{Calculation and Logic:} \\
When $n=0$, $H=1$. This suggests a constant thickness added to the stack regardless of the number of books.

\vspace{15pt}
Answer: \colorbox{yellow!100}{C}
\end{solution}

\vspace{20pt}

% Question 31
% Category: Practice | Difficulty: Hard | Question Type: Multiple Choice
\question
A car rental company charges \$45 per day. Another company charges \$25 per day plus \$0.25 per mile. If a person rents a car for one day, at what mileage $m$ will the costs be equal?
\begin{tasks}(4)
    \task 80 miles
    \task 100 miles
    \task 180 miles
    \task 280 miles
\end{tasks}

\begin{solution}
\textbf{Conceptual Explanation:} \\
Set the two cost equations equal to each other.

\vspace{10pt}
\textbf{Calculation and Logic:} \\
$45 = 25 + 0.25m$ \\
$20 = 0.25m$ \\
$m = 20 / 0.25 = 80$.

\vspace{15pt}
Answer: \colorbox{yellow!100}{A}
\end{solution}

\vspace{20pt}

% Question 32
% Category: Practice | Difficulty: Medium | Question Type: Multiple Choice
\question
An ice block weighs 100 lbs and melts at a rate of 4 lbs per hour. Which function represents the weight $W$ after $h$ hours?
\begin{tasks}(2)
    \task $W = 100 + 4h$
    \task $W = 4h - 100$
    \task $W = 100 - 4h$
    \task $W = 96h$
\end{tasks}

\begin{solution}
\textbf{Conceptual Explanation:} \\
Melting implies a decrease in weight over time (negative slope).

\vspace{10pt}
\textbf{Calculation and Logic:} \\
Initial = 100. Rate = $-4$. \\
Equation: $W = 100 - 4h$.

\vspace{15pt}
Answer: \colorbox{yellow!100}{C}
\end{solution}

\vspace{20pt}

% Question 33
% Category: Practice | Difficulty: Hard | Question Type: Multiple Choice
\question
A balloon is being inflated. Its volume $V$, in cubic inches, $t$ seconds after inflation starts is $V = 15t + 10$. If the maximum volume is 400 cubic inches, what is the maximum value of $t$?
\begin{tasks}(4)
    \task 26
    \task 27.3
    \task 26.6
    \task 25
\end{tasks}

\begin{solution}
\textbf{Conceptual Explanation:} \\
Set the volume to the maximum limit and solve for $t$.

\vspace{10pt}
\textbf{Calculation and Logic:} \\
$400 = 15t + 10$ \\
$390 = 15t$ \\
$t = 26$.

\vspace{15pt}
Answer: \colorbox{yellow!100}{A}
\end{solution}

\vspace{20pt}

% Question 34
% Category: Practice | Difficulty: Medium | Question Type: Multiple Choice
\question
A store offers a discount where the price $P$ of a jacket is $P = 120 - 5w$, where $w$ is the number of weeks it has been on the shelf. What is the original price of the jacket?
\begin{tasks}(4)
    \task \$5
    \task \$115
    \task \$120
    \task \$24
\end{tasks}

\begin{solution}
\textbf{Conceptual Explanation:} \\
Original price corresponds to $w=0$.

\vspace{10pt}
\textbf{Calculation and Logic:} \\
$P = 120 - 5(0) = 120$.

\vspace{15pt}
Answer: \colorbox{yellow!100}{C}
\end{solution}

\vspace{20pt}

% Question 35
% Category: Practice | Difficulty: Hard | Question Type: Multiple Choice
\question
A pump can remove water from a basement at a rate of 12 gallons per minute. If the basement starts with 1,500 gallons and the pump runs for 2 hours, how many gallons of water remain?
\begin{tasks}(4)
    \task 1,476
    \task 780
    \task 60
    \task 1,440
\end{tasks}

\begin{solution}
\textbf{Conceptual Explanation:} \\
Convert hours to minutes, calculate the total removed, and subtract from the initial amount.

\vspace{10pt}
\textbf{Calculation and Logic:} \\
2 hours = 120 minutes. \\
Removed = $12 \times 120 = 1,440$ gallons. \\
Remaining = $1,500 - 1,440 = 60$.

\vspace{15pt}
Answer: \colorbox{yellow!100}{C}
\end{solution}

\vspace{20pt}

% Question 36
% Category: Practice | Difficulty: Medium | Question Type: Multiple Choice
\question
A line graph shows that for every 3 units increase in $x$, $y$ increases by 12 units. If the line passes through $(0, 5)$, what is the linear model for $y$?
\begin{tasks}(2)
    \task $y = 3x + 5$
    \task $y = 4x + 5$
    \task $y = 12x + 5$
    \task $y = \frac{1}{4}x + 5$
\end{tasks}

\begin{solution}
\textbf{Conceptual Explanation:} \\
Calculate the slope ($m$) first, then use the given $y$-intercept.

\vspace{10pt}
\textbf{Calculation and Logic:} \\
$m = 12 / 3 = 4$. \\
$b = 5$. \\
Equation: $y = 4x + 5$.

\vspace{15pt}
Answer: \colorbox{yellow!100}{B}
\end{solution}

\vspace{20pt}

% Question 37
% Category: Practice | Difficulty: Hard | Question Type: Multiple Choice
\question
The total cost $C$ for a manufacturer to produce $n$ items is $C = 5n + 2,000$. If the manufacturer sells each item for \$15, how many items must be produced and sold to break even (where cost equals revenue)?
\begin{tasks}(4)
    \task 100
    \task 200
    \task 133
    \task 150
\end{tasks}

\begin{solution}
\textbf{Conceptual Explanation:} \\
Set Cost equal to Revenue ($15n$).

\vspace{10pt}
\textbf{Calculation and Logic:} \\
$15n = 5n + 2,000$ \\
$10n = 2,000$ \\
$n = 200$.

\vspace{15pt}
Answer: \colorbox{yellow!100}{B}
\end{solution}

\vspace{20pt}

% Question 38
% Category: Practice | Difficulty: Medium | Question Type: Multiple Choice
\question

The graph above shows the weight $W$, in pounds, of a puppy over $x$ weeks. If the puppy gained 2 pounds every week and started at 8 pounds, which equation matches the graph?
\begin{tasks}(2)
    \task $W = 8x + 2$
    \task $W = 2x + 8$
    \task $W = 2x + 10$
    \task $W = 10x + 8$
\end{tasks}

\begin{solution}
\textbf{Conceptual Explanation:} \\
"Started at" is the $b$ value. "Gained every week" is the $m$ value.

\vspace{10pt}
\textbf{Calculation and Logic:} \\
$m = 2, b = 8$. \\
Equation: $W = 2x + 8$.

\vspace{15pt}
Answer: \colorbox{yellow!100}{B}
\end{solution}

\vspace{20pt}

% Question 39
% Category: Practice | Difficulty: Hard | Question Type: Multiple Choice
\question
A hiker descends from an elevation of 5,000 feet at a rate of 250 feet per hour. Which function gives the elevation $E$ after $h$ hours?
\begin{tasks}(2)
    \task $E = 250h - 5,000$
    \task $E = 5,000 - 250h$
    \task $E = 5,000 + 250h$
    \task $E = 5,000h - 250$
\end{tasks}

\begin{solution}
\textbf{Conceptual Explanation:} \\
Initial elevation is high and decreases over time.

\vspace{10pt}
\textbf{Calculation and Logic:} \\
Initial = 5,000. Rate = $-250$. \\
Equation: $E = 5,000 - 250h$.

\vspace{15pt}
Answer: \colorbox{yellow!100}{B}
\end{solution}

\vspace{20pt}

% Question 40
% Category: Practice | Difficulty: Medium | Question Type: Multiple Choice
\question
A subscription box service charges \$20 for the first box and \$15 for each subsequent box. Which equation gives the total cost $C$ for $n$ boxes?
\begin{tasks}(2)
    \task $C = 15n + 20$
    \task $C = 15(n - 1) + 20$
    \task $C = 20n + 15$
    \task $C = 35n$
\end{tasks}

\begin{solution}
\textbf{Conceptual Explanation:} \\
The \$20 only applies once. The \$15 applies to the remaining ($n-1$) boxes.

\vspace{10pt}
\textbf{Calculation and Logic:} \\
$C = 20 + 15(n - 1)$. \\
(This simplifies to $15n + 5$, but Choice B is the correct setup).

\vspace{15pt}
Answer: \colorbox{yellow!100}{B}
\end{solution}

\vspace{20pt}

% Question 41
% Category: Practice | Difficulty: Hard | Question Type: Multiple Choice
\question
The amount of snow on the ground $S$ is modeled by $S = 2t + 4$ during a storm. If the storm lasts from 10:00 AM to 4:00 PM, what is the total increase in snow?
\begin{tasks}(4)
    \task 4 inches
    \task 6 inches
    \task 12 inches
    \task 16 inches
\end{tasks}

\begin{solution}
\textbf{Conceptual Explanation:} \\
Total increase is the rate multiplied by the total duration.

\vspace{10pt}
\textbf{Calculation and Logic:} \\
Duration: 10 AM to 4 PM = 6 hours. \\
Rate = 2 inches per hour. \\
Increase = $2 \times 6 = 12$ inches.

\vspace{15pt}
Answer: \colorbox{yellow!100}{C}
\end{solution}

\vspace{20pt}

% Question 42
% Category: Practice | Difficulty: Medium | Question Type: Multiple Choice
\question
A scientist finds that the temperature $T$ of a sample drops 3 degrees every minute. If the starting temperature was 75 degrees, what is the temperature after 10 minutes?
\begin{tasks}(4)
    \task 72 degrees
    \task 45 degrees
    \task 30 degrees
    \task 105 degrees
\end{tasks}

\begin{solution}
\textbf{Conceptual Explanation:} \\
Calculate the total drop and subtract it from the initial temperature.

\vspace{10pt}
\textbf{Calculation and Logic:} \\
Drop = $3 \times 10 = 30$. \\
$75 - 30 = 45$.

\vspace{15pt}
Answer: \colorbox{yellow!100}{B}
\end{solution}

\vspace{20pt}

% Question 43
% Category: Practice | Difficulty: Hard | Question Type: Multiple Choice
\question
A company's revenue $R$ is modeled by $R = 100x$ and its costs $C$ are $C = 40x + 1,200$. What is the profit equation ($P = R - C$)?
\begin{tasks}(2)
    \task $P = 60x + 1,200$
    \task $P = 60x - 1,200$
    \task $P = 140x + 1,200$
    \task $P = 140x - 1,200$
\end{tasks}

\begin{solution}
\textbf{Conceptual Explanation:} \\
Subtract the entire cost expression from the revenue expression.

\vspace{10pt}
\textbf{Calculation and Logic:} \\
$P = 100x - (40x + 1,200)$ \\
$P = 100x - 40x - 1,200$ \\
$P = 60x - 1,200$.

\vspace{15pt}
Answer: \colorbox{yellow!100}{B}
\end{solution}

\vspace{20pt}

% Question 44
% Category: Practice | Difficulty: Medium | Question Type: Multiple Choice
\question
A car fuel tank holds 15 gallons. The car uses 0.04 gallons per mile. Which equation gives the gallons $G$ remaining after $m$ miles?
\begin{tasks}(2)
    \task $G = 0.04m - 15$
    \task $G = 15 + 0.04m$
    \task $G = 15 - 0.04m$
    \task $G = 15.04m$
\end{tasks}

\begin{solution}
\textbf{Conceptual Explanation:} \\
Fuel decreases with distance.

\vspace{10pt}
\textbf{Calculation and Logic:} \\
Initial = 15. Rate = $-0.04$. \\
Equation: $G = 15 - 0.04m$.

\vspace{15pt}
Answer: \colorbox{yellow!100}{C}
\end{solution}

\vspace{20pt}

% Question 45
% Category: Practice | Difficulty: Hard | Question Type: Multiple Choice
\question

Line A has a growth rate of 10 units per day. Line B has a growth rate of 5 units per day but starts with 50 units. If Line A starts with 0 units, on what day will Line A exceed Line B?
\begin{tasks}(4)
    \task Day 5
    \task Day 10
    \task Day 11
    \task Day 20
\end{tasks}

\begin{solution}
\textbf{Conceptual Explanation:} \\
Set the inequality $A > B$.

\vspace{10pt}
\textbf{Calculation and Logic:} \\
Line A: $10d$. Line B: $5d + 50$. \\
$10d > 5d + 50$ \\
$5d > 50$ \\
$d > 10$. \\
The first day it exceeds is Day 11.

\vspace{15pt}
Answer: \colorbox{yellow!100}{C}
\end{solution}

\vspace{20pt}

% Question 46
% Category: Practice | Difficulty: Medium | Question Type: Multiple Choice
\question
A plumber charges \$50 for the first hour and \$35 for each additional hour. Which model represents the total cost $C$ for $h$ hours ($h \ge 1$)?
\begin{tasks}(2)
    \task $C = 35h + 50$
    \task $C = 35(h - 1) + 50$
    \task $C = 50h + 35$
    \task $C = 85h$
\end{tasks}

\begin{solution}
\textbf{Conceptual Explanation:} \\
The \$50 is only for the first hour. Subsequent hours are charged at the lower rate.

\vspace{10pt}
\textbf{Calculation and Logic:} \\
$C = 50 + 35(h - 1)$.

\vspace{15pt}
Answer: \colorbox{yellow!100}{B}
\end{solution}

\vspace{20pt}

% Question 47
% Category: Practice | Difficulty: Hard | Question Type: Multiple Choice
\question
If the value of an asset $V$ is modeled by $V = 50,000 - 4,000t$, how many years will it take for the asset to lose half of its original value?
\begin{tasks}(4)
    \task 6.25 years
    \task 12.5 years
    \task 5 years
    \task 10 years
\end{tasks}

\begin{solution}
\textbf{Conceptual Explanation:} \\
Original value is 50,000. Half is 25,000. Solve for $t$.

\vspace{10pt}
\textbf{Calculation and Logic:} \\
$25,000 = 50,000 - 4,000t$ \\
$-25,000 = -4,000t$ \\
$t = 6.25$.

\vspace{15pt}
Answer: \colorbox{yellow!100}{A}
\end{solution}

\vspace{20pt}

% Question 48
% Category: Practice | Difficulty: Medium | Question Type: Multiple Choice
\question
A gardener plants a tree that is 4 feet tall. It grows 0.5 feet per year. How tall will it be after 14 years?
\begin{tasks}(4)
    \task 7 feet
    \task 11 feet
    \task 18 feet
    \task 9 feet
\end{tasks}

\begin{solution}
\textbf{Conceptual Explanation:} \\
Total growth = rate $\times$ time. Add this to the starting height.

\vspace{10pt}
\textbf{Calculation and Logic:} \\
Growth = $0.5 \times 14 = 7$. \\
$4 + 7 = 11$.

\vspace{15pt}
Answer: \colorbox{yellow!100}{B}
\end{solution}

\vspace{20pt}

% Question 49
% Category: Practice | Difficulty: Hard | Question Type: Multiple Choice
\question
A data plan charges a flat fee of \$40 and \$10 for every GB over 5 GB. If a customer's bill was \$90, how many total GB did they use?
\begin{tasks}(4)
    \task 5 GB
    \task 9 GB
    \task 10 GB
    \task 14 GB
\end{tasks}

\begin{solution}
\textbf{Conceptual Explanation:} \\
Determine how much was paid in overage fees, calculate the overage GB, and add to the base 5 GB.

\vspace{10pt}
\textbf{Calculation and Logic:} \\
Overage fees = $90 - 40 = 50$. \\
Overage GB = $50 / 10 = 5$. \\
Total GB = $5 (\text{base}) + 5 (\text{overage}) = 10$.

\vspace{15pt}
Answer: \colorbox{yellow!100}{C}
\end{solution}

\vspace{20pt}

% Question 50
% Category: Practice | Difficulty: Medium | Question Type: Multiple Choice
\question
An airplane descends from 35,000 feet to 5,000 feet in 15 minutes. What is its constant rate of descent in feet per minute?
\begin{tasks}(4)
    \task 1,000
    \task 2,000
    \task 3,000
    \task 2,333
\end{tasks}

\begin{solution}
\textbf{Conceptual Explanation:} \\
Rate is total change divided by total time.

\vspace{10pt}
\textbf{Calculation and Logic:} \\
Change = $35,000 - 5,000 = 30,000$ feet. \\
Time = 15 minutes. \\
Rate = $30,000 / 15 = 2,000$.

\vspace{15pt}
Answer: \colorbox{yellow!100}{B}
\end{solution}









% =============================================================
% Question End
% =============================================================

\end{questions}
\begin{center}
\rule{.5\textwidth}{1pt}
\end{center}
\end{document}
